\vspace*{-20cm} 
\chapter{Introducción:}
\section{Motivación: }

En el ámbito del laboratorio Medical Analytics (en adelante MEDAL) en el Centro de Tecnología Biomédica (en adelante CTB) al estar ligado a la Universidad Politécnica de Madrid (en adelante UPM) se trata de un laboratorio con alta rotación de personal. En la rotación de personal se pueden distinguir los siguientes roles: profesores, doctores, estudiantes a doctor, contratados y estudiantes de grado o máster. Cada nuevo integrante dependiendo de sus necesidades necesita acceso a diferentes servicios heterogéneos como pueden ser redes de internet, acceso a servidores, canales de comunicación entre otros. 

\notaseguridad{La gestión manual de estos privilegios incrementa el vector de ataque, pudiendo dejar cuentas activas tras la desvinculación de los usuarios.}



Un factor crítico identificado en el entorno del laboratorio es la gestión del acceso a los servidores. Actualmente, el acceso se realiza de forma nativa y compartida, existiendo una gran cantidad de usuarios con privilegios de administrador (\textit{root} o \textit{sudoers}). 

Esta falta de granularidad en el control de acceso ha derivado en problemas técnicos. Al carecer de entornos aislados, las instalaciones de librerías y controladores por parte de distintos investigadores han generado conflictos de dependencias y corrupciones en el sistema operativo, comprometiendo la estabilidad de los servicios y la continuidad de las investigaciones en curso.

\notaseguridad{La cantidad de usuarios con permisos de administrador se contrapone al \textbf{Principio de Mínimo Privilegio} aumentando el riesgo de errores accidentelas que pueden inutilizar la infraestructura de los servidores.}




\section{Objetivos del proyecto:}
Este Trabajo Fin de Grado surge de la necesidad de centralizar y automatizar el ciclo de vida de los usuarios dentro del grupo de investigación MEDAL. El objetivo principal es desarrollar una plataforma web que actúe como panel de control (dashboard), permitiendo la administración y provisión automática de servicios. 

Esta solución no solo busca optimizar la carga administrativa, sino también mitigar los problemas derivados del \textbf{acceso nativo persistente} a los servidores del laboratorio. La plataforma propuesta permitirá orquestar estos accesos de forma granular, asegurando que la configuración de los servicios se mantenga íntegra y que cada investigador cuente con un entorno de trabajo estable y seguro, alineado con las necesidades específicas de su proyecto de investigación.

\section{Estructura de la memoria:}
Para la estructuración del proyecto el documento se va a organizar de la siguiente manera:
\begin{itemize}
  TODO: Poner los incrementos.

\end{itemize}
